\documentclass[10pt,a4paper]{amsart}
\usepackage[margin=19mm]{geometry}
\usepackage{amsmath,amssymb,mathtools}
\usepackage[scr=rsfs,cal=euler]{mathalfa}
\usepackage{xfrac}
\usepackage[unicode,hidelinks,bookmarks=false]{hyperref}
\newcommand{\ie}{i.e.\ }
\newcommand{\cf}{cf.\ }
\newcommand{\resp}{respectively\ }
\newcommand{\Subsection}{\subsection}
\providecommand{\nobreakdash}{\nobreak\hbox{-}}
\setlength{\emergencystretch}{2em}
\multlinegap0pt
\usepackage{amssymb}
\usepackage{float}
\newtheorem{proposition}{Proposition}
\newtheorem{lemma}{Lemma}
\newcommand\bbR{\mathbb R}
\newcommand\cC{\mathcal{C}}
\newcommand\cF{\mathcal F}
\newcommand\pa{\partial}
\newcommand\tw{\tilde w_{0+}}
\title{Solving the Scattering Problem for Open Wave-Guide Networks, I\\Fundamental Solutions and Integral Equations}
\author{Charles L. Epstein}
\date{Source: arXiv:2302.04353v4 (CC BY 4.0)}
\begin{document}
\maketitle

\noindent\emph{Source and licence.} Solving the Scattering Problem for Open Wave-Guide Networks, I\\Fundamental Solutions and Integral Equations. Version arXiv:2302.04353v4 (CC BY 4.0), distributed by arXiv under the Creative Commons Attribution 4.0 International licence, http://creativecommons.org/licenses/by/4.0/, as recorded in the arXiv licence metadata for this exact version and shown on the abstract page https://arxiv.org/abs/2302.04353v4. Full licence notice: \texttt{licenses/arxiv-2302.04353-CC-BY-4.0.txt}.\par\smallskip

\noindent\textit{Editorial scope.} Counted unit: the article's proposition prop2.94 (source0.tex lines 2118--2130). The article's complete original proof of it is delivered with it (source0.tex lines 2131--2279, one \texttt{proof} environment of 149 lines). No mathematics is written, repaired or re-derived: the statement and the proof are the article's own lines, in the article's own notation, and no retained line is substituted or re-flowed.  Retained uncounted as the source-local context the proof needs: lines 1081--1083; lines 3277--3279; lines 3452--3454.  Nothing was omitted from any retained span, so the package carries no omission record.  No retained line contains a citation, so the wrapper carries no bibliography and no bibliography entry.  No retained line contains a figure environment, an image-inclusion command or drawing code, so no image is delivered and the package declares no asset dependency.  Editorial formatting only, in addition: an AMS \texttt{amsart} preamble that restates the article's own declarations and macros used by the retained lines, namely the environments \texttt{proposition}, \texttt{lemma} on separate counters, together with the standard packages those lines need.  The retained environments print in this document as Proposition 1, Lemma 1 in order of appearance, whereas the article prints the same items with the numbering of its own full document.  The complete native source of the article as distributed in the arXiv e-print for this exact version is bundled unchanged as \texttt{sources/137145/source0.tex}.
\begin{equation}\label{eqn26.51}
  \cF_{x_1}[(i/4)H^{(1)}_0(k_1|x-y|)](\xi)=\frac{ie^{i|x_2-y_2|\sqrt{k_1^2-\xi^2}}e^{-iy_1\xi}}{2\sqrt{k_1^2-\xi^2}};
\end{equation}

    \begin{equation}\label{eqn73.9}
      |\tw(\xi,x_2;0,y_2)|\leq M\frac{e^{-\sqrt{\xi^2-k_1^2}(x_2+y_2-2d)}}{(1+|\xi|)^{3}}.
    \end{equation}

   \begin{equation}\label{eqn92.9}
     |\tw(\xi,x_2;0,y_2)|<M\frac{e^{-\sqrt{\xi^2-k_2^2}|x_2-y_2|}}{(1+|\xi|)^2}\left[\frac{1}{(1+|\xi|)}+|y_2-x_2|\right],
   \end{equation}

\begin{proposition}\label{prop2.94}
  Let $(\sigma,\tau)\in \cC_{\alpha}(\bbR)\oplus\cC_{\alpha+\frac 12}(\bbR)$ and
  let $v_{n}^{l,r}(x_2)e^{i\xi_n^{l,r}x_1}$ be a wave-guide mode for
  $\Delta+k_1^2+q_{l,r}(x_2),$ then, for all  $\pm x_1>0,$ 
  \begin{equation}
    \begin{split}
     &\int\limits_{-\infty}^{\infty}\int\limits_{-\infty}^{\infty}[g_{k_1}(|(x_1,x_2-y_2)|)+
        w^{c;l,r}_{0+}(x_1,x_2;0,y_2)]\tau(y_2)v_n^{l,r}(x_2)dy_2dx_2=0,\\
     &\int\limits_{-\infty}^{\infty}\int\limits_{-\infty}^{\infty}\pa_{y_1}[g_{k_1}(|(x_1-y_1,x_2-y_2)|)+\\
        &\phantom{mmmmmmmmmmmm}w^{c;l,r}_{0+}(x_1,x_2;y_1,y_2)]_{y_1=0}\sigma(y_2)v_n^{l,r}(x_2)dy_2dx_2=0.
    \end{split}
  \end{equation}
\end{proposition}

\begin{proof}
  We give the details for the right half plane. The estimates proved in the
  previous sections show that these integrals are absolutely convergent and
  therefore we can change the order of the integrations. To prove the
  proposition we show that, for each $y_2$ and $x_1>0,$ we have
   \begin{equation}\label{eqn212.93} 
    \begin{split}
     &\int\limits_{-\infty}^{\infty}[g_{k_1}(|(x_1,x_2-y_2)|)+
        w^{c;r}_{0+}(x_1,x_2;0,y_2)]v_n^{r}(x_2)dx_2=0,\\
     &\int\limits_{-\infty}^{\infty}\pa_{y_1}[g_{k_1}(|(x_1-y_1,x_2-y_2)|)+
        w^{c;r}_{0+}(x_1,x_2;y_1,y_2)]_{y_1=0}v_n^{r}(x_2)dx_2=0.
    \end{split}
  \end{equation}
  To prove these statements we use the Sommerfeld integral representation for
  the free space fundamental solution, see~\eqref{eqn26.51}.

  We begin with the single layer term, which can be written as the iterated
  integral:
  \begin{equation}\label{eqn213.91}
    \frac{1}{2\pi}
    \int\limits_{-\infty}^{\infty}\left[\int\limits_{-\infty}^{\infty}\frac{ie^{i|x_2-y_2|\sqrt{k_1^2-\xi^2}}}{2\sqrt{k_1^2-\xi^2}}+
       \int\limits_{\Gamma_{\nu}^+} \tw^{c;r}(\xi,x_2;y_2)\right] e^{ix_1\xi }d\xi\,v_n^{r}(x_2)dx_2.
  \end{equation}
  We would like to change the order of integrations in this integral, which
  would be  easily justified if the integral in $\xi$ were over any finite
  interval. Note that
  \begin{equation}
    \int\limits_{k_1+1}^{\infty}\left|\frac{ie^{-|x_2-y_2|\sqrt{\xi^2-k_1^2}}e^{-iy_1\xi}}{2\sqrt{\xi^2-k_1^2}}\right|d\xi\leq
    M[1+|\log|x_2-y_2|].
  \end{equation}
  Combining this with the estimates~\eqref{eqn73.9} and~\eqref{eqn92.9}, and the fact that
  $|v_n^r(x_2)|\leq Me^{-|x_2|\sqrt{\xi_n^{r\,2}-k_1^2}}$ shows that these integrals are
  absolutely convergent and therefore we can interchange the order of the
  integrations. By analyticity we can also replace the integral in the first term with an
  integral over $\Gamma_{\nu}^+,$ to obtain
   \begin{equation}
        \frac{1}{2\pi}
        \int\limits_{\Gamma_{\nu}^+}\left[\int\limits_{-\infty}^{\infty}
          \frac{ie^{i|x_2-y_2|\sqrt{k_1^2-\xi^2}}}{2\sqrt{k_1^2-\xi^2}}+
       \tw^{c;r}(\xi,x_2;y_2)\right] e^{ix_1\xi}v_n^{r}(x_2)dx_2d\xi.
   \end{equation}
   Using the fact that $\xi^{r\, 2}_nv_n^r=(\pa_{x_2}^2+k_1^2+q_r(x_2))v_n^r$ and
   integrating by parts  it follows that
   \begin{equation}\label{eqn216.93}
     \int\limits_{-\infty}^{\infty}\left[\frac{ie^{i|x_2-y_2|\sqrt{k_1^2-\xi^2}}}{2\sqrt{k_1^2-\xi^2}}+
       \tw^{c;r}(\xi,x_2;y_2)\right] e^{ix_1\xi}v_n^{r}(x_2)dx_2=\frac{ie^{ix_1\xi}v_n^r(y_2)}{\xi_n^{r\,2}-\xi^2}.
   \end{equation}
   Hence the double integral is
   \begin{equation}
     \frac{1}{2\pi i}\int\limits_{\Gamma_{\nu}^+}\frac{ie^{ix_1\xi}v_n^r(y_2)d\xi}{\xi^2-\xi_n^{r\,2}},
   \end{equation}
   which, using Cauchy's theorem, is easily seen to vanish for $x_1>0.$

   The double layer is almost the same; the double integral in~\eqref{eqn213.91} is
   replaced by
     \begin{equation}\label{eqn218.93}
    \frac{i}{2\pi}
    \int\limits_{-\infty}^{\infty}\left[\,\int\limits_{\Gamma_{\nu}^+}
    \left[\frac{ie^{i|x_2-y_2|\sqrt{k_1^2-\xi^2}}}{2\sqrt{k_1^2-\xi^2}}+
       \tw^{c;r}(\xi,x_2;y_2)\right]\xi e^{ix_1\xi}d\xi\right] v_n^{r}(x_2)dx_2.
     \end{equation}
     The integral involving $\tw^{c;r}(\xi,x_2;y_2)$ is again easily seen to be
     absolutely convergent, but the additional factor of $\xi$ makes the other
     term more subtle. We need an estimate for this term that takes account of
     the fact that $x_1>0.$
     \begin{lemma}
       For $x_1>0, \lambda>0$ and $R>k_1$ let
       \begin{equation}
       f^{\pm}_R(x_1,\lambda)=\pm\int\limits_{\pm R}^{\pm \infty}
       \frac{e^{ix_1\xi-\lambda\sqrt{\xi^2-k_1^2}}\xi d\xi}{\sqrt{\xi^2-k_1^2}}.
       \end{equation}
       These functions satisfy the estimates
       \begin{equation}
         |f^{\pm}_R(x_1,\lambda)|\leq \frac{Me^{-\lambda\sqrt{R^2-k_1^2}}}{x_1}
       \end{equation}
     \end{lemma}
     \begin{proof}
       As $f^-_R(x_1,\lambda)=\overline{f^+_R(x_1,\lambda)}$ it suffices to do the $+$-case.
       If we let $s=\sqrt{\xi^2-k_1},$ then
       \begin{equation}
         f^+_R(x_1,\lambda)=\int\limits_{\sqrt{R^2-k_1^2}}^{\infty}e^{ix_1\sqrt{s^2+k_1^2}-\lambda s}ds.
       \end{equation}
       Noting that
       \begin{multline}\label{eqn222.93}
         \pa_s\left[\frac{\sqrt{s^2+k_1^2}}{s}e^{ix_1\sqrt{s^2+k_1^2}}\right]=\\
         ix_1e^{ix_1\sqrt{s^2+k_1^2}}-\frac{k_1^2}{s^2\sqrt{s^2+k_1^2}}e^{ix_1\sqrt{s^2+k_1^2}},
       \end{multline}
       integration by parts shows that
       \begin{multline}
         f^+_R(x_1,\lambda)=
         \frac{1}{ix_1}\Bigg[
           \frac{k_1e^{-\lambda\sqrt{R^2-k_1^2}+ix_1R}}{\sqrt{R^2-k_1^2}}
           +\\
           \int\limits_{\sqrt{R^2-k_1^2}}^{\infty}\left[\frac{\lambda\sqrt{s^2+k_1^2}}{s}+
             \frac{k_1^2}{s^2\sqrt{s^2+k_1^2}}\right]e^{ix_1\sqrt{s^2+k_1^2}-\lambda
             s}ds\Bigg].
       \end{multline}
      The integral is easily seen to be $O(e^{-\lambda\sqrt{R^2-k_1^2}}),$ which
      completes the proof of the lemma.
     \end{proof}

     We rewrite the integral in~\eqref{eqn218.93} as
     \begin{multline}
       \frac{i}{2\pi}
    \int\limits_{-\infty}^{\infty}\int\limits_{\Gamma_{\nu}^+}
    \left[\frac{ie^{i|x_2-y_2|\sqrt{k_1^2-\xi^2}}}{2\sqrt{k_1^2-\xi^2}}\right]\xi
    e^{ix_1\xi}d\xi\, v_n^{r}(x_2)dx_2=\\
    \frac{i}{2\pi}
    \int\limits_{-\infty}^{\infty}\left[\,\int\limits_{\Gamma_{\nu}^+\cap D_R}
      \left[\frac{ie^{i|x_2-y_2|\sqrt{k_1^2-\xi^2}}}{2\sqrt{k_1^2-\xi^2}}\right]\xi
      e^{ix_1\xi}d\xi+f_R^+(x_1,|x_2-y_2|)+
    f_R^-(x_1,|x_2-y_2|)\right] v_n^{r}(x_2)dx_2.
     \end{multline}
     In the part of the integral over $\Gamma_{\nu}^+\cap D_R$ we can
     interchange the order of the integrations. We use the lemma to estimate
     the other terms
     \begin{equation}
       \begin{split}
       \Bigg| \int\limits_{-\infty}^{\infty}[f_R^+(x_1,|x_2-y_2|)&+
         f_R^-(x_1,|x_2-y_2|)] v_n^{r}(x_2)dx_2\Bigg| \\
       &\leq
       \frac{M}{x_1}\int_{-\infty}^{\infty}e^{-\sqrt{R^2-k_1^2}|x_2-y_2|}e^{-\sqrt{\xi_n^{r\,2}-k_1^2}|x_2|}dx_2\\
       &\leq\frac{M}{x_1(\sqrt{R^2-k_1^2}-\sqrt{\xi_n^{r\,2}-k_1^2})}.
       \end{split}
     \end{equation}
     Thus we see that
       \begin{multline}\label{eqn227.93}
       \int\limits_{-\infty}^{\infty}\left[\,\int\limits_{\Gamma_{\nu}^+}
    \left[\frac{ie^{i|x_2-y_2|\sqrt{k_1^2-\xi^2}}}{2\sqrt{k_1^2-\xi^2}}+
       \tw^{c;r}(\xi,x_2;y_2)\right]\xi e^{ix_1\xi}d\xi\right]
    v_n^{r}(x_2)dx_2=\\
    \lim_{R\to\infty}
     \int\limits_{\Gamma_{\nu}^+\cap  D_R}\int\limits_{-\infty}^{\infty}\left[
    \left[\frac{ie^{i|x_2-y_2|\sqrt{k_1^2-\xi^2}}}{2\sqrt{k_1^2-\xi^2}}+
       \tw^{c;r}(\xi,x_2;y_2)\right]v_n^{r}(x_2)dx_2\right]\,\xi e^{ix_1\xi}
    d\xi.
       \end{multline}
       The $x_2$-integral is computed in~\eqref{eqn216.93}, giving
     \begin{multline}\label{eqn228.93}
      \int\limits_{-\infty}^{\infty}\left[\,\int\limits_{\Gamma_{\nu}^+}
    \left[\frac{ie^{i|x_2-y_2|\sqrt{k_1^2-\xi^2}}}{2\sqrt{k_1^2-\xi^2}}+
       \tw^{c;r}(\xi,x_2;y_2)\right]\xi e^{ix_1\xi}d\xi\right]
    v_n^{r}(x_2)dx_2=\\
    \lim_{R\to\infty}
     v_n^{r}(y_2)\int\limits_{\Gamma_{\nu}^+\cap  D_R}\frac{i\xi e^{ix_1\xi}}{\xi_n^{r\,2}-\xi^2} d\xi=0.   
     \end{multline}
    The last equality follows from Cauchy's theorem. It completes the proof
    of~\eqref{eqn212.93} and also of the proposition.
\end{proof}

\end{document}
